\chapter{Temporal Networks: When Time Changes Everything}
\label{chap:temporal_networks}

\textit{This chapter builds the mathematical language of temporal networks from the ground up. We
start with a simple three-person story that makes the core problem immediately tangible, then
formalise every concept step by step. By the end you will be able to read temporal network
papers, compute reachability by hand, and understand exactly how the project codebase stores
and uses contact sequences.}

%=============================================================================
\section{Why Static Graphs Fail for Epidemics: The Phantom Path Problem}
\label{sec:phantom_paths}
%=============================================================================

\subsection{A Three-Person Story}

Imagine three people: Alice (A), Bob (B), and Carol (C). They meet as follows:

\begin{itemize}
    \item \textbf{Monday ($t = 1$):} Alice meets Bob. They shake hands.
    \item \textbf{Tuesday ($t = 2$):} Bob meets Carol. They share a meal.
    \item \textbf{Sunday ($t = 7$):} Carol meets Alice. They sit together at a concert.
\end{itemize}

Now suppose Bob was infected with a virus on Monday. The question an epidemiologist asks is:
\textit{who could Bob have infected, and through whom?}

\begin{examplebox}{Drawing the Static Graph}
If we draw a graph that simply records ``who has met whom'' (ignoring time), we get:

\medskip
\begin{center}
\begin{tikzpicture}[
  nd/.style={circle, draw=blue!60!black, fill=blue!8, thick,
             minimum size=10mm, font=\bfseries\small},
  >=Stealth, thick
]
  \node[nd] (A) at (0,0)   {A};
  \node[nd] (B) at (2,0)   {B};
  \node[nd] (C) at (1,1.7) {C};
  \draw (A) -- node[below, font=\scriptsize]{Mon} (B);
  \draw (B) -- node[right, font=\scriptsize]{Tue} (C);
  \draw (C) -- node[left,  font=\scriptsize]{Sun} (A);
\end{tikzpicture}
\end{center}

\medskip
In this static graph every pair of nodes is connected. A standard graph algorithm will
happily report the path $C \to A \to B$: Carol can infect Alice, and Alice can infect Bob.

But wait. Carol meets Alice on \textbf{Sunday} ($t=7$). Alice meets Bob on \textbf{Monday}
($t=1$). Sunday comes \textit{after} Monday. A pathogen cannot travel backwards in time. Carol
cannot infect Alice on Sunday and then have that infection arrive at Bob on Monday, six days
earlier.

The path $C \to A \to B$ \textbf{does not exist in reality}. It is a \textbf{phantom path}:
a path that looks valid in the static graph but is causally impossible because the edge
timestamps are in the wrong order.
\end{examplebox}

\begin{conceptbox}{Phantom Path}
A \textbf{phantom path} is a sequence of edges in the static projection of a temporal network
that appears to be a valid path but violates causal ordering: at least one edge in the sequence
has a timestamp that is \textit{earlier} than a preceding edge. An epidemic cannot travel
along a phantom path.
\end{conceptbox}

\begin{keyinsight}{The Core Problem with Static Graphs}
Static graphs tell you \textit{who has been in contact}, but not \textit{in what order}. For
epidemic spreading, order is everything: a virus can only move \emph{forward} in time. When a
static-graph algorithm suggests that Carol is a likely source because she is ``connected'' to
both Alice and Bob, it may be completely wrong---if Carol's contacts all happened after the
outbreak was already established, she cannot be the origin.

\textbf{This is why temporal networks exist}: they preserve the one piece of information that
static graphs discard---the timestamps that encode causality.
\end{keyinsight}

\subsection{A Brief History of Temporal Network Analysis}

Temporal or \textit{dynamic} graph models have been studied since at least the 1970s in the
context of communications networks, but the field gained coherent identity around 2012 with the
comprehensive review by \cite{holme2012temporal}. The key empirical observation motivating this
work was that real-world contact sequences exhibit strong non-Poissonian (``bursty'') timing:
inter-event times follow heavy-tailed rather than exponential distributions, with dramatic
consequences for how fast pathogens spread.

The formal study of \textit{time-respecting paths} and \textit{temporal reachability} was
pioneered by Kempe, Kleinberg, and Kumar \cite{kempe2000connectivity}, who established
complexity results for temporal connectivity problems. The event graph representation used in
this thesis derives from Saram\"{a}ki, Kivel\"{a}, and Karsai \cite{saramaki2019temporal},
covered in depth in Chapter~8.

%=============================================================================
\section{Formal Definition of a Temporal Network}
\label{sec:tn_definition}
%=============================================================================

We now give a precise mathematical definition. Every symbol will be introduced before it is
used.

\begin{conceptbox}{Notation We Need Before the Definition}
\begin{itemize}[noitemsep]
  \item $\mathcal{V}$ --- the \textbf{node set}: the collection of all individuals (people,
        computers, animals) in the network. We write $|\mathcal{V}| = n$ for the number of
        nodes.
  \item $\times$ --- the \textbf{Cartesian product}: $\mathcal{V} \times \mathcal{V}$ is the
        set of all \emph{ordered pairs} of nodes, e.g.\ $\{(A,B), (A,C), (B,A), \ldots\}$.
  \item $\mathcal{T}$ --- the \textbf{observation window}: the time interval $[t_{\min},
        t_{\max}]$ during which we recorded contacts.
  \item $\mathbb{R}^+$ --- the set of \textbf{non-negative real numbers}: $\{x \in \mathbb{R}
        \mid x \geq 0\}$.
  \item $\subseteq$ --- ``is a subset of'': every element on the left also belongs to the set
        on the right.
  \item $(u, v, t)$ --- an ordered \textbf{triple}: the contact between node $u$ and node $v$
        at time $t$.
  \item $\in$ --- ``is an element of'': $e \in E^T$ means event $e$ belongs to the event set.
\end{itemize}
\end{conceptbox}

\begin{conceptbox}{Temporal Network: Formal Definition}
A \textbf{temporal network} (also called a \emph{contact sequence}) is a triple
\[
  G^T = \bigl(\mathcal{V},\; E^T,\; \mathcal{T}\bigr)
\]
where:
\begin{itemize}[noitemsep]
  \item $\mathcal{V} = \{v_1, v_2, \ldots, v_n\}$ is the finite set of $n$ \textbf{nodes}.
  \item $\mathcal{T} = [t_{\min}, t_{\max}] \subset \mathbb{R}^+$ is the
        \textbf{observation window}.
  \item $E^T \subseteq \mathcal{V} \times \mathcal{V} \times \mathcal{T}$ is the
        finite set of \textbf{temporal events} (contacts).
\end{itemize}
Each event $e = (u, v, t) \in E^T$ records that nodes $u$ and $v$ were in contact at time $t$.
We assume contacts are \textbf{undirected}: $(u, v, t) \in E^T$ if and only if
$(v, u, t) \in E^T$ (if Alice meets Bob, Bob also meets Alice). The total number of distinct
events is $|E^T|$.
\end{conceptbox}

\textbf{Reading the definition aloud}: $G^T = (\mathcal{V}, E^T, \mathcal{T})$ is read as
``$G$-superscript-$T$ is a temporal network with node set $\mathcal{V}$, event set $E^T$, and
observation window $\mathcal{T}$.'' The superscript $T$ stands for \emph{temporal} (not
transpose).

\begin{examplebox}{The Alice--Bob--Carol Network as a Formal Temporal Network}
Returning to our three-person story:
\[
  \mathcal{V} = \{A,\; B,\; C\}, \qquad
  \mathcal{T} = [1, 7], \qquad
  E^T = \{(A,B,1),\; (B,C,2),\; (C,A,7)\}
\]
Reading each triple:
\begin{itemize}[noitemsep]
  \item $(A, B, 1)$: Alice and Bob were in contact at time $t = 1$ (Monday).
  \item $(B, C, 2)$: Bob and Carol were in contact at time $t = 2$ (Tuesday).
  \item $(C, A, 7)$: Carol and Alice were in contact at time $t = 7$ (Sunday).
\end{itemize}

As a table (the most common format in data files):
\begin{center}
\begin{tabular}{ccc}
\toprule
\textbf{Node $u$} & \textbf{Node $v$} & \textbf{Time $t$} \\
\midrule
A & B & 1 \\
B & C & 2 \\
C & A & 7 \\
\bottomrule
\end{tabular}
\end{center}

The static projection $G = (\mathcal{V}, \mathcal{E})$ discards the timestamps and retains
only connectivity:
$\mathcal{E} = \{\{A,B\}, \{B,C\}, \{C,A\}\}$.
This is a triangle. Every node can reach every other node. But only \textit{two} of the three
directed paths are causally valid, as we will see in the next section.
\end{examplebox}

\begin{analogybox}{The Photograph vs.\ the Video}
A static graph is a \textbf{photograph}: it shows who was in the same frame, but gives no
information about who arrived first or who left before whom. A temporal network is a
\textbf{video}: it shows every frame in order. To trace an epidemic, you need the video. The
photograph might show Alice and Bob standing next to a sick person, but only the video reveals
whether they were there before or after the person fell ill.
\end{analogybox}

\subsection{Distinguishing Temporal Networks from Related Concepts}

The literature uses several overlapping terms. Here is a precise glossary:

\begin{conceptbox}{Terminology: Contact Sequence vs.\ Snapshot vs.\ Dynamic Graph}
\begin{itemize}[noitemsep]
  \item \textbf{Contact sequence / temporal network:} An explicit list of timestamped events
        $(u, v, t)$. Time is treated as a fine-grained attribute. This is the format used in
        this thesis and by the TSIR simulation engine.
  \item \textbf{Time-varying graph / snapshot sequence:} A sequence of static graphs
        $\{G_1, G_2, \ldots, G_T\}$, each valid over a time bin
        $[t_{k-1}, t_k)$. This is a discretised approximation; within-bin ordering is lost.
  \item \textbf{Dynamic graph:} A broad term for any graph that changes over time, including
        models where edges are permanently added or deleted (e.g.\ social network growth).
  \item \textbf{Multilayer / supra-adjacency:} A representation stacking all snapshots into a
        single block matrix, connecting copies of each node across time layers.
\end{itemize}
The key structural difference: a contact sequence preserves the \textit{exact timing} of every
event, whereas a snapshot sequence only preserves whether a contact occurred within a time bin.
\end{conceptbox}

%=============================================================================
\section{Time-Respecting Paths: The Causal Highways of Epidemics}
\label{sec:time_respecting_paths}
%=============================================================================

\subsection{The Relay Race Analogy}

\begin{analogybox}{The Relay Race}
Imagine a relay race. Runner A passes the baton to Runner B at the $t = 10\,\text{s}$ mark.
Runner B then passes the baton to Runner C at the $t = 25\,\text{s}$ mark.

This is valid: B \textit{received} the baton at $t = 10\,\text{s}$ and \textit{passed} it at
$t = 25\,\text{s}$. The sequence is $10 < 25$: time moves forward.

Now imagine the race announcement says B passed the baton to C at $t = 8\,\text{s}$. This is
impossible: B had not yet received the baton at $t = 8\,\text{s}$. B cannot pass what B does
not yet have.

An epidemic works exactly the same way. A node can only \textit{transmit} a pathogen after it
has \textit{received} it. Every step in the chain must happen \emph{after} the previous step.
This is the defining constraint of a \textbf{time-respecting path}.
\end{analogybox}

\subsection{Formal Definition}

\begin{mathbox}{Time-Respecting Path: Step-by-Step Definition}
Let $G^T = (\mathcal{V}, E^T, \mathcal{T})$ be a temporal network.

\textbf{Step 1 --- What is a path?} A \textit{path} from node $v_0$ to node $v_k$ is a
sequence of nodes and events:
\[
  P = \bigl(v_0,\; e_1,\; v_1,\; e_2,\; v_2,\; \ldots,\; e_k,\; v_k\bigr)
\]
where each $e_i = (v_{i-1},\, v_i,\, t_i) \in E^T$ is a contact event involving consecutive
nodes in the path.

\textbf{Step 2 --- What makes it ``time-respecting''?} The path is \textbf{time-respecting}
if and only if the event times are \textit{strictly increasing}:
\[
  t_1 < t_2 < \cdots < t_k.
\]
This enforces causality: each node must receive the pathogen before it can pass it on.

\textbf{Step 3 --- What is the temporal length?} The \textbf{length} of $P$ is $k$, the
number of hops. The \textbf{duration} of $P$ is $t_k - t_1$, the elapsed time from first to
last event.

\textbf{Step 4 --- What does it mean for $u$ to reach $v$?} Node $u$ \textbf{temporally
reaches} node $v$ (written $u \leadsto v$) if there exists \textit{at least one}
time-respecting path from $u$ to $v$.
\end{mathbox}

\subsection{Worked Example: Five Events, Valid and Invalid Paths}

We now work through a complete example. Consider the following five-event contact sequence on
four nodes $\{1, 2, 3, 4\}$:

\begin{examplebox}{Five-Event Temporal Network}
\begin{center}
\begin{tabular}{cccc}
\toprule
\textbf{Event} & \textbf{Node $u$} & \textbf{Node $v$} & \textbf{Time $t$} \\
\midrule
$e_1$ & 1 & 2 & 1 \\
$e_2$ & 2 & 3 & 3 \\
$e_3$ & 3 & 4 & 5 \\
$e_4$ & 1 & 3 & 4 \\
$e_5$ & 4 & 2 & 2 \\
\bottomrule
\end{tabular}
\end{center}

\medskip
\noindent\textbf{Which paths from node 1 are time-respecting?}

We check every possible sequence of events that starts at node 1:

\begin{center}
\begin{tabular}{llll}
\toprule
\textbf{Candidate path} & \textbf{Event times} & \textbf{Strictly increasing?} & \textbf{Valid?} \\
\midrule
$1 \xrightarrow{e_1} 2$               & $t=1$             & --- (single hop)  & \textcolor{green!60!black}{\textbf{Yes}} \\
$1 \xrightarrow{e_4} 3$               & $t=4$             & --- (single hop)  & \textcolor{green!60!black}{\textbf{Yes}} \\
$1 \xrightarrow{e_1} 2 \xrightarrow{e_2} 3$  & $t=1 < t=3$  & Yes  & \textcolor{green!60!black}{\textbf{Yes}} \\
$1 \xrightarrow{e_1} 2 \xrightarrow{e_2} 3 \xrightarrow{e_3} 4$ & $t=1<t=3<t=5$ & Yes & \textcolor{green!60!black}{\textbf{Yes}} \\
$1 \xrightarrow{e_4} 3 \xrightarrow{e_3} 4$ & $t=4 < t=5$ & Yes & \textcolor{green!60!black}{\textbf{Yes}} \\
$1 \xrightarrow{e_1} 2 \xrightarrow{e_5?} 4$ & need $e$ with nodes $(2,4)$ after $t=1$: none & --- & \textcolor{red!70!black}{\textbf{No path}} \\
$1 \xrightarrow{e_4} 3 \xrightarrow{e_2?} 2$ & $e_2=(2,3,t=3)$, but $t=3 < t=4$ & $3 \not> 4$ & \textcolor{red!70!black}{\textbf{Invalid}} \\
$1 \xrightarrow{e_1} 2 \xrightarrow{e_5} 4$ & $e_5=(4,2,t=2)$; $t=2>t=1$ but $4\to 2$ direction wrong via $e_5$: need $(2,\cdot)$ event after $t=1$; $e_5$ is $(4,2,t=2)$, so 2 is reached \textit{by} 4, not \textit{from} 2 & n/a & \textcolor{red!70!black}{\textbf{No}} \\
\bottomrule
\end{tabular}
\end{center}

\medskip
\textbf{Temporal reachability set of node 1:} $\mathcal{I}(1) = \{1, 2, 3, 4\}$ --- node 1
can reach all others.

\textbf{Can node 4 reach node 1?} Node 4 appears in events $e_3$ (at $t=5$, with node 3) and
$e_5$ (at $t=2$, with node 2). From $e_5$: $4 \to 2$ at $t=2$. From node 2: the only later
event is $e_2$ ($2$--$3$ at $t=3$). From node 3: $e_3$ ($3$--$4$ at $t=5$). So $4 \leadsto 2
\leadsto 3 \leadsto 4$ --- but this loops back and never reaches node 1. There is no path from
4 to 1. \textbf{Asymmetry confirmed}: $1 \leadsto 4$ but $4 \not\leadsto 1$.
\end{examplebox}

%=============================================================================
\section{The $\Delta t$ Constraint: Modelling the Infectious Period}
\label{sec:delta_t_constraint}
%=============================================================================

\subsection{Why We Need an Additional Constraint}

A simple time-respecting path allows any gap between consecutive events, even a gap of years.
But in a real epidemic, a person who was infected does not remain infectious forever: they
either recover or die. If node $B$ was infected at time $t=1$ and the next event involving $B$
is at time $t=100$, can $B$ really transmit the virus at $t=100$?

Not if the virus only survives in the host for, say, 10 days. This motivates the
$\Delta t$-constrained path, where $\Delta t$ (read ``delta $t$'') is the maximum time a node
can remain infectious.

\begin{conceptbox}{The $\Delta t$ Constraint, Symbol by Symbol}
\[
  t_i - t_{i-1} \leq \Delta t \quad \forall\, i \in \{2, \ldots, k\}
\]
Reading this symbol by symbol:
\begin{itemize}[noitemsep]
  \item $t_i$ --- the time of the $i$-th event in the path.
  \item $t_{i-1}$ --- the time of the immediately preceding event.
  \item $t_i - t_{i-1}$ --- the \textbf{gap} (inter-event time) between consecutive hops.
  \item $\leq$ --- ``less than or equal to''.
  \item $\Delta t$ --- the \textbf{maximum infectious duration}: the longest a node can hold
        the pathogen without transmitting it before recovering.
  \item $\forall\, i \in \{2, \ldots, k\}$ --- ``for every hop index $i$ from 2 to $k$'': the
        constraint must hold at \textit{every} step, not just the first or last.
\end{itemize}
\end{conceptbox}

\begin{mathbox}{$\Delta t$-Constrained Time-Respecting Path}
A time-respecting path
\[
  P = (v_0,\; e_1,\; v_1,\; e_2,\; v_2,\; \ldots,\; e_k,\; v_k)
\]
with event times $t_1 < t_2 < \cdots < t_k$ is \textbf{$\Delta t$-constrained} if and only
if:
\[
  t_i - t_{i-1} \leq \Delta t \quad \text{for all } i = 2, 3, \ldots, k.
\]
The set of $\Delta t$-constrained paths is always a subset of the set of all time-respecting
paths. As $\Delta t \to \infty$, the constraint disappears and we recover unconstrained
time-respecting paths.

In our TSIR simulations (Chapter~3), $\Delta t \approx 1/\mu$ in expectation, where $\mu$ is
the SIR recovery rate.
\end{mathbox}

\subsection{Numerical Example: Cutting Paths with $\Delta t = 2$}

\begin{examplebox}{Which Paths Survive a $\Delta t = 2$ Cut?}
We return to our five-event network. The valid time-respecting paths from node 1 were:
\begin{enumerate}[noitemsep]
  \item $1 \xrightarrow{t=1} 2$: single hop, no gap to check. \textbf{Survives} any $\Delta t$.
  \item $1 \xrightarrow{t=4} 3$: single hop. \textbf{Survives}.
  \item $1 \xrightarrow{t=1} 2 \xrightarrow{t=3} 3$: gap = $3 - 1 = 2 \leq 2$. \textbf{Survives}.
  \item $1 \xrightarrow{t=1} 2 \xrightarrow{t=3} 3 \xrightarrow{t=5} 4$: gaps $= 2$ and $= 2$.
        Both $\leq 2$. \textbf{Survives}.
  \item $1 \xrightarrow{t=4} 3 \xrightarrow{t=5} 4$: gap $= 5 - 4 = 1 \leq 2$. \textbf{Survives}.
\end{enumerate}

Now try $\Delta t = 1$:
\begin{enumerate}[noitemsep]
  \item $1 \to 2$: single hop. \textbf{Survives}.
  \item $1 \to 3$: single hop. \textbf{Survives}.
  \item $1 \xrightarrow{1} 2 \xrightarrow{3} 3$: gap $= 3 - 1 = 2 > 1$. \textbf{Cut!}
  \item $1 \xrightarrow{1} 2 \xrightarrow{3} 3 \xrightarrow{5} 4$: first gap already $=2>1$.
        \textbf{Cut!}
  \item $1 \xrightarrow{4} 3 \xrightarrow{5} 4$: gap $= 1 \leq 1$. \textbf{Survives}.
\end{enumerate}

With $\Delta t = 1$, node 1 can no longer reach node 2 via 3, and the long chain $1 \to 2 \to
3 \to 4$ is severed. The reachability set from node 1 shrinks:
\[
  \Delta t = \infty:\; \mathcal{I}(1) = \{1,2,3,4\}, \qquad
  \Delta t = 2:\; \mathcal{I}(1) = \{1,2,3,4\}, \qquad
  \Delta t = 1:\; \mathcal{I}(1) = \{1,2,3,4\}
\]
(In this small example the set happens to be the same, but the \textit{number of paths} and
the \textit{routes} available differ. In larger, sparser networks the set itself shrinks
significantly.)
\end{examplebox}

\begin{keyinsight}{Why $\Delta t$ Matters for Source Detection}
When we search for the epidemic source, we ask: ``from which node could the observed infected
set be reached by a valid transmission chain?'' If we forget the $\Delta t$ constraint, we
over-estimate the set of possible sources (any node connected in the static graph looks
plausible). If we set $\Delta t$ too small, we under-estimate it (we incorrectly rule out the
true source). Using the correct $\Delta t$ --- matching the biology of the disease --- gives
the tightest, most informative set of candidate sources.
\end{keyinsight}

%=============================================================================
\section{Temporal vs.\ Static Reachability: The Numbers That Differ}
\label{sec:reachability_comparison}
%=============================================================================

In this section we compute both static and temporal reachability for a small, fully worked
example and display them in a comparison table. The numbers will disagree, and every
disagreement corresponds to a phantom path in the static graph.

\subsection{A Four-Node, Five-Event Example}

Consider the following contact sequence on nodes $\{A, B, C, D\}$:

\begin{center}
\begin{tabular}{cccc}
\toprule
\textbf{Event} & \textbf{$u$} & \textbf{$v$} & \textbf{$t$} \\
\midrule
$e_1$ & A & B & 1 \\
$e_2$ & B & C & 2 \\
$e_3$ & C & D & 3 \\
$e_4$ & D & A & 5 \\
$e_5$ & A & C & 4 \\
\bottomrule
\end{tabular}
\end{center}

\begin{figure}[H]
\centering
\begin{tikzpicture}[
  nd/.style={circle, draw=blue!60!black, fill=blue!8, thick,
             minimum size=10mm, font=\bfseries},
  ev/.style={font=\scriptsize, color=blue!60!black},
  >=Stealth, thick
]
% Nodes in a square
\node[nd] (A) at (0, 2)  {A};
\node[nd] (B) at (2, 4)  {B};
\node[nd] (C) at (4, 2)  {C};
\node[nd] (D) at (2, 0)  {D};

% Edges with timestamps
\draw (A) -- node[ev, left]{$e_1,\,t=1$} (B);
\draw (B) -- node[ev, right]{$e_2,\,t=2$} (C);
\draw (C) -- node[ev, right]{$e_3,\,t=3$} (D);
\draw (D) -- node[ev, left]{$e_4,\,t=5$} (A);
\draw (A) -- node[ev, above]{$e_5,\,t=4$} (C);
\end{tikzpicture}
\caption{The four-node, five-event temporal network used in the reachability comparison.
Edge labels show the event identifier and its timestamp.}
\label{fig:reachability_example}
\end{figure}

\subsection{Computing Static Reachability}

In the static projection $G = (\{A,B,C,D\}, \{\{A,B\},\{B,C\},\{C,D\},\{D,A\},\{A,C\}\})$,
we simply ask: is there any path (ignoring time) from $u$ to $v$?

Since the graph is connected (you can traverse from any node to any other), the answer is
\textbf{yes for all pairs}. The static reachability matrix $R^{\text{static}}$ where
$R^{\text{static}}_{uv} = 1$ iff $u$ can reach $v$ in $G$:

\[
R^{\text{static}} =
\begin{array}{c|cccc}
  & A & B & C & D \\
\hline
A & 1 & 1 & 1 & 1 \\
B & 1 & 1 & 1 & 1 \\
C & 1 & 1 & 1 & 1 \\
D & 1 & 1 & 1 & 1 \\
\end{array}
\]

All entries are 1: every node reaches every other node. (We include diagonal entries as 1
by convention: every node trivially reaches itself.)

\subsection{Computing Temporal Reachability}

Now we apply the strict time-ordering rule. We trace all time-respecting paths from each node.

\begin{examplebox}{Temporal Reachability: Node A}
Events involving A: $e_1$ (A--B, $t=1$), $e_5$ (A--C, $t=4$), $e_4$ (D--A, $t=5$, but this
is \textit{incoming} to A).

\textbf{From A, outgoing in time:}
\begin{itemize}[noitemsep]
  \item $e_1$: $A \xrightarrow{t=1} B$. From B: $e_2$ (B--C, $t=2>1$) $\Rightarrow B \xrightarrow{t=2} C$. From C: $e_3$ (C--D, $t=3>2$) $\Rightarrow C \xrightarrow{t=3} D$. Also $e_5$ (A--C, $t=4>2$) but we came \textit{to} C, not from A again.
  \item $e_5$: $A \xrightarrow{t=4} C$. From C: $e_3$ (C--D, $t=3$). But $3 < 4$, so this is \textbf{not} valid. No forward events from C after $t=4$.
\end{itemize}
$\mathcal{I}(A) = \{A, B, C, D\}$ (all reachable).
\end{examplebox}

\begin{examplebox}{Temporal Reachability: Node B}
Events involving B: $e_1$ (A--B, $t=1$, incoming), $e_2$ (B--C, $t=2$, outgoing).

\textbf{From B outgoing:}
$B \xrightarrow{t=2} C \xrightarrow{t=3} D$. Node D: event $e_4$ (D--A, $t=5>3$), so
$D \xrightarrow{t=5} A$. But this is an event \textit{involving} D and A; since we arrived at
D at $t=3$, and $e_4$ is at $t=5>3$, A is reachable.

Can B reach A? Only through D at $t=5$. Yes: $B \to C \to D \to A$ with times $2<3<5$.

$\mathcal{I}(B) = \{B, C, D, A\}$.
\end{examplebox}

\begin{examplebox}{Temporal Reachability: Node C}
Events involving C: $e_2$ (B--C, $t=2$), $e_3$ (C--D, $t=3$), $e_5$ (A--C, $t=4$).

\textbf{From C outgoing:}
\begin{itemize}[noitemsep]
  \item $C \xrightarrow{t=3} D$ (via $e_3$). From D: $e_4$ (D--A, $t=5>3$). So $C \to D \to A$.
  \item Event $e_5$ is at $t=4>3$ but it involves A and C. Since we are \textit{at} C after
        time $t=3$, we need events where C is the \textit{source}. $e_5$ is the contact $(A,C,4)$:
        this is A meeting C at $t=4$. If we are already at C, can we reach A via $e_5$? Yes: $C
        \xrightarrow{e_5,\,t=4} A$ (undirected contact, so C can transmit to A at $t=4$ as long
        as $4 > 3$). (\textit{Wait}: we need to check: does this make a valid path? $C$ starts at
        time $0$ (or whenever), and is involved in $e_3$ at $t=3$ and $e_5$ at $t=4$. The path
        $C \xrightarrow{t=3} D$ leads to D, not a return to C.)
\end{itemize}

Let us be systematic. From node C at the \textit{start} of the observation:
\begin{itemize}[noitemsep]
  \item Direct 1-hop: $C \xrightarrow{t=3} D$; $C \xrightarrow{t=4} A$ (via $e_5$, undirected).
  \item From D at $t=3$: $D \xrightarrow{t=5} A$ (via $e_4$). Conflict: we also reached A via $e_5$ at $t=4$.
  \item From A at $t=4$ (reached via $e_5$): no outgoing events from A after $t=4$ except $e_4$ which is \textit{incoming} to A.
\end{itemize}

Can C reach B? B only appears in $e_1$ ($t=1$) and $e_2$ ($t=2$), both \textit{before} C's
first event at $t=3$. So B is \textbf{not reachable} from C via time-respecting paths.

$\mathcal{I}(C) = \{C, D, A\}$. Note: \textbf{B is not reachable from C}!
\end{examplebox}

\begin{examplebox}{Temporal Reachability: Node D}
Events involving D: $e_3$ (C--D, $t=3$, incoming), $e_4$ (D--A, $t=5$).

\textbf{From D outgoing:} $D \xrightarrow{t=5} A$ (via $e_4$). From A at $t=5$: no later
events. From A, can we reach B or C after $t=5$? No events involve A, B, or C after $t=5$.

$\mathcal{I}(D) = \{D, A\}$. D can only reach itself and A.
\end{examplebox}

\subsection{Side-by-Side Comparison}

\begin{examplebox}{Static vs.\ Temporal Reachability Matrix}

\begin{center}
\begin{tabular}{cc}
\begin{tabular}{c|cccc}
\multicolumn{5}{c}{\textbf{Static reachability} $R^{\text{static}}$}\\[2pt]
  & A & B & C & D \\
\hline
A & 1 & 1 & 1 & 1 \\
B & 1 & 1 & 1 & 1 \\
C & 1 & 1 & 1 & 1 \\
D & 1 & 1 & 1 & 1 \\
\end{tabular}
&
\begin{tabular}{c|cccc}
\multicolumn{5}{c}{\textbf{Temporal reachability} $R^T$}\\[2pt]
  & A & B & C & D \\
\hline
A & 1 & 1 & 1 & 1 \\
B & 1 & 1 & 1 & 1 \\
C & 1 & \textcolor{red!70!black}{\textbf{0}} & 1 & 1 \\
D & 1 & \textcolor{red!70!black}{\textbf{0}} & \textcolor{red!70!black}{\textbf{0}} & 1 \\
\end{tabular}
\end{tabular}
\end{center}

\medskip
The entries marked in \textcolor{red!70!black}{\textbf{red}} are \textbf{phantom paths}: the
static graph claims $C \leadsto B$ and $D \leadsto B$ and $D \leadsto C$, but no
time-respecting path exists for these pairs.

\textbf{Explanation of each phantom:}
\begin{itemize}[noitemsep]
  \item $C \to B$: B's only events are at $t=1$ and $t=2$. C's earliest outgoing event is at
        $t=3$. No valid path from C can reach B at a time after C was first involved.
  \item $D \to B$: D's earliest outgoing event is at $t=5$. B's last event is at $t=2$. No
        path from D can travel back to $t=2$ to reach B.
  \item $D \to C$: D can only reach A (at $t=5$). From A there are no later events involving
        C.
\end{itemize}

Note also the \textbf{asymmetry}: $R^T_{AB} = 1$ (A reaches B) but $R^T_{BA} = 1$ too in
this example because the cycle A--B--C--D--A eventually loops. In general, temporal
reachability need not be symmetric even in undirected networks.
\end{examplebox}

%=============================================================================
\section{Adjacency Representation of Temporal Networks}
\label{sec:adjacency_representation}
%=============================================================================

\subsection{The Adjacency Tensor $\mathbf{A}(t)$}

For a static graph we use the \textbf{adjacency matrix} $A \in \{0,1\}^{n \times n}$, where
$A_{ij} = 1$ iff there is an edge between $i$ and $j$. For a temporal network, we need to
capture the time dimension as well.

\begin{conceptbox}{The Adjacency Tensor}
Define $A_{ij}(t) = 1$ if nodes $i$ and $j$ were in contact at time $t$, and $A_{ij}(t) = 0$
otherwise. This gives a function $\mathbf{A} : \mathcal{T} \to \{0,1\}^{n \times n}$, which
we can think of as a \textbf{three-dimensional tensor} (a cube of numbers) with dimensions:
\[
  \underbrace{n}_{\text{node }i} \times \underbrace{n}_{\text{node }j} \times
  \underbrace{T}_{\text{time slice }t}
\]
For each time step $t$, the slice $\mathbf{A}(\cdot, \cdot, t) = A_t$ is the adjacency matrix
of the \textbf{snapshot graph} at time $t$.

For \textbf{instantaneous contacts} (our case), $A_{ij}(t) = 1$ only at isolated time points.
For \textbf{discrete time}, $t$ takes integer values and we can write $\mathbf{A}$ as a stack
of $T$ binary matrices of size $n \times n$.
\end{conceptbox}

\begin{examplebox}{Adjacency Tensor for the Four-Node Example}
Using the events $\{(A,B,1), (B,C,2), (C,D,3), (A,C,4), (D,A,5)\}$ and discrete time steps
$t \in \{1,2,3,4,5\}$:

\medskip
\noindent$A_1$ (contacts at $t=1$): only $A$--$B$.
\[
A_1 = \begin{pmatrix} 0&1&0&0\\ 1&0&0&0\\ 0&0&0&0\\ 0&0&0&0 \end{pmatrix}
\]

\noindent$A_2$ (contacts at $t=2$): only $B$--$C$.
\[
A_2 = \begin{pmatrix} 0&0&0&0\\ 0&0&1&0\\ 0&1&0&0\\ 0&0&0&0 \end{pmatrix}
\]
And so on for $t=3,4,5$. At all other time steps (and all steps not listed), $A_t = \mathbf{0}$.
\end{examplebox}

\subsection{The COO Format: How Computers Store Sparse Temporal Graphs}

The adjacency tensor is very sparse: most entries are zero (most pairs of nodes are not in
contact at any given moment). Storing the full $n \times n \times T$ tensor wastes enormous
amounts of memory. Instead, computers use the \textbf{coordinate (COO) format}: store only the
non-zero entries as a list of (row, column, value) triples.

For a temporal network, the COO format stores each event $(u, v, t)$ explicitly. This is
exactly the event list we have been using all along.

\begin{conceptbox}{COO Format for Temporal Networks}
The COO event list for a temporal network with $|E^T|$ events is:
\[
\underbrace{\begin{pmatrix} u_1 \\ v_1 \\ t_1 \end{pmatrix},\;
\begin{pmatrix} u_2 \\ v_2 \\ t_2 \end{pmatrix},\;
\ldots,\;
\begin{pmatrix} u_{|E^T|} \\ v_{|E^T|} \\ t_{|E^T|} \end{pmatrix}}_{\text{one column per event}}
\]
This is stored as a matrix with 3 rows and $|E^T|$ columns.

In \textbf{PyTorch Geometric} (the GNN library used in this project), the graph connectivity is
stored in \texttt{edge\_index}: a $2 \times |E|$ tensor where the first row contains source
node indices and the second row contains target node indices. Edge features (including
timestamps) are stored separately in \texttt{edge\_attr}.
\end{conceptbox}

\begin{examplebox}{From Event List to PyTorch Geometric Format}
Our four-node example has events (using integer node labels $0=A, 1=B, 2=C, 3=D$):
\begin{itemize}[noitemsep]
  \item $(A,B,1) \to (0,1,1)$
  \item $(B,C,2) \to (1,2,2)$
  \item $(C,D,3) \to (2,3,3)$
  \item $(A,C,4) \to (0,2,4)$
  \item $(D,A,5) \to (3,0,5)$
\end{itemize}

The PyTorch Geometric \texttt{edge\_index} for the static projection (both directions, since
undirected):
\[
\texttt{edge\_index} = \begin{pmatrix}
0 & 1 & 1 & 2 & 2 & 3 & 0 & 3 \\
1 & 0 & 2 & 1 & 3 & 2 & 2 & 0
\end{pmatrix}
\]
(first row: sources; second row: targets; both directions included for each undirected edge).

The timestamps become \texttt{edge\_attr}:
\[
\texttt{edge\_attr} = [1,\; 1,\; 2,\; 2,\; 3,\; 3,\; 4,\; 5]
\]
corresponding to the edges in the same column order. This is exactly the format read by our
\texttt{gnn/static\_gnn.py} and temporal GNN models.
\end{examplebox}

%=============================================================================
\section{Temporal Centrality Metrics}
\label{sec:temporal_centrality}
%=============================================================================

Classical centrality metrics like degree, betweenness, and closeness have \textit{temporal}
counterparts that carry substantially more epidemiological information. We define each one,
give intuition, and then compute it step by step on our running example.

\subsection{Temporal Degree}

\begin{conceptbox}{Temporal Degree: Definition}
The \textbf{temporal degree} of node $v$ is the total number of events (contacts) involving
$v$ in the observation window:
\[
  k_v^T = \bigl|\{(u, v, t) \in E^T\}\bigr| + \bigl|\{(v, w, t) \in E^T\}\bigr|
\]
In plain English: count every event where $v$ appears on either side.

This differs from the \textbf{static degree} $k_v = |\mathcal{N}(v)|$, which counts the
number of \textit{distinct neighbours} of $v$. A node with many repeated contacts to the same
neighbour has high temporal degree but the same static degree as a node with one contact to
that neighbour.
\end{conceptbox}

\begin{examplebox}{Computing Temporal Degree for All Four Nodes}
Event list: $\{(A,B,1), (B,C,2), (C,D,3), (A,C,4), (D,A,5)\}$.

\begin{center}
\begin{tabular}{lll}
\toprule
\textbf{Node} & \textbf{Events involving this node} & \textbf{Temporal degree $k_v^T$} \\
\midrule
A & $(A,B,1)$, $(A,C,4)$, $(D,A,5)$ & 3 \\
B & $(A,B,1)$, $(B,C,2)$             & 2 \\
C & $(B,C,2)$, $(C,D,3)$, $(A,C,4)$ & 3 \\
D & $(C,D,3)$, $(D,A,5)$             & 2 \\
\bottomrule
\end{tabular}
\end{center}

Static degrees: A has 3 neighbours ($B,C,D$), B has 2, C has 3, D has 2. In this particular
example the numbers match because no node has repeated contacts with the same neighbour.

\textbf{What can temporal degree tell us?} High temporal degree means a node is frequently
involved in contacts. In an epidemic context, a node with many events is likely to be infected
early and to transmit widely. But temporal degree alone does not capture \textit{when} those
events occur or whether they can form valid transmission chains.
\end{examplebox}

\subsection{Temporal Betweenness Centrality}

\begin{conceptbox}{Temporal Betweenness: Definition}
The \textbf{temporal betweenness centrality} of node $v$ measures how often $v$ lies on the
\textit{fastest} (shortest hop count) time-respecting paths between all other pairs of nodes:
\[
  C_B^T(v) = \sum_{\substack{s \neq v \neq d \\ s,d \in \mathcal{V}}}
             \frac{\sigma^T(s, d \mid v)}{\sigma^T(s, d)}
\]
where:
\begin{itemize}[noitemsep]
  \item $\sigma^T(s, d)$ --- the number of shortest (minimum-hop) time-respecting paths from
        $s$ to $d$.
  \item $\sigma^T(s, d \mid v)$ --- the number of those shortest paths that pass through $v$.
  \item The fraction $\sigma^T(s, d \mid v) / \sigma^T(s, d)$ is the proportion of shortest
        paths through $v$ for the pair $(s, d)$.
\end{itemize}
Summing over all ordered pairs $(s, d)$ with $s \neq v \neq d$ gives the total centrality.
\end{conceptbox}

\begin{examplebox}{Temporal Betweenness: Step-by-Step for Four Nodes}
We use strict time-ordering and unit edge weights (hop count). We already know the
time-respecting paths, so we identify the \textit{shortest} (fewest hops) ones for each pair.

\medskip
\noindent\textbf{All ordered pairs and their shortest time-respecting paths:}

\begin{center}
\begin{tabular}{llll}
\toprule
$(s,d)$ & Shortest temporal path(s) & Hops & Intermediaries \\
\midrule
$(A,B)$ & $A \xrightarrow{1} B$ & 1 & none \\
$(A,C)$ & $A \xrightarrow{4} C$ and $A \xrightarrow{1} B \xrightarrow{2} C$ & 1 and 2 & shortest: none (direct hop) \\
$(A,D)$ & $A \xrightarrow{1} B \xrightarrow{2} C \xrightarrow{3} D$ and $A \xrightarrow{4} C \xrightarrow{?}$ (no later C--D event after $t=4$) & 3 & B, C \\
$(B,C)$ & $B \xrightarrow{2} C$ & 1 & none \\
$(B,D)$ & $B \xrightarrow{2} C \xrightarrow{3} D$ & 2 & C \\
$(B,A)$ & $B \xrightarrow{2} C \xrightarrow{3} D \xrightarrow{5} A$ & 3 & C, D \\
$(C,D)$ & $C \xrightarrow{3} D$ & 1 & none \\
$(C,A)$ & $C \xrightarrow{4} A$ and $C \xrightarrow{3} D \xrightarrow{5} A$ & 1 and 2 & shortest: none (direct hop) \\
$(D,A)$ & $D \xrightarrow{5} A$ & 1 & none \\
$(C,B)$ & \textbf{none} (phantom) & --- & --- \\
$(D,B)$ & \textbf{none} (phantom) & --- & --- \\
$(D,C)$ & \textbf{none} (phantom) & --- & --- \\
\bottomrule
\end{tabular}
\end{center}

\medskip
\noindent\textbf{Counting betweenness contributions:}

For each node $v$, sum over all pairs $(s,d)$ the fraction of shortest temporal paths through
$v$:

\begin{itemize}
  \item \textbf{Node B:}
    \begin{itemize}[noitemsep]
      \item $(A,D)$: shortest path is $A \to B \to C \to D$ (3 hops). B is on this path.
            Contribution: $1/1 = 1$.
      \item $(B,A)$: B is the source, not an intermediary. Not counted.
      \item No other shortest paths pass through B.
    \end{itemize}
    $C_B^T(B) = 1$.

  \item \textbf{Node C:}
    \begin{itemize}[noitemsep]
      \item $(A,D)$: $A \to B \to C \to D$. C is on this path. Contribution: $1/1 = 1$.
      \item $(B,D)$: $B \to C \to D$. C is on this path. Contribution: $1/1 = 1$.
      \item $(B,A)$: $B \to C \to D \to A$. C is on this path. Contribution: $1/1 = 1$.
    \end{itemize}
    $C_B^T(C) = 3$.

  \item \textbf{Node D:}
    \begin{itemize}[noitemsep]
      \item $(B,A)$: $B \to C \to D \to A$. D is on this path. Contribution: $1/1 = 1$.
    \end{itemize}
    $C_B^T(D) = 1$.

  \item \textbf{Node A:} We check all pairs where A is neither source nor destination:
    \begin{itemize}[noitemsep]
      \item $(B,C)$: $B \to C$ directly. A not involved.
      \item $(B,D)$: $B \to C \to D$. A not involved.
      \item $(C,D)$: $C \to D$ directly. A not involved.
      \item $(D,A)$: D is source, A is destination.
    \end{itemize}
    $C_B^T(A) = 0$.
\end{itemize}

\medskip
\noindent\textbf{Summary:}
\begin{center}
\begin{tabular}{cc}
\toprule
\textbf{Node} & \textbf{Temporal betweenness $C_B^T(v)$} \\
\midrule
A & 0 \\
B & 1 \\
C & \textbf{3} \\
D & 1 \\
\bottomrule
\end{tabular}
\end{center}
\end{examplebox}

\begin{keyinsight}{Static vs.\ Temporal Centrality Can Disagree Strongly}
In the static graph, all nodes have the same degree (A: 3, B: 2, C: 3, D: 2) and the static
betweenness of each node is influenced by the number of shortest static paths passing through
it. In this connected, roughly symmetric graph, static betweenness would distribute roughly
evenly.

The temporal betweenness, however, concentrates on node \textbf{C}: C is on 3 out of 4 pairs
that have temporal paths through an intermediary. This is because C is in the ``middle'' of
the time ordering: it is reachable early (at $t=2$ from B, at $t=3$ from the chain starting
at A) and has outgoing temporal contacts later.

\textbf{Epidemiological consequence:} If we were looking for the most important superspreader
in this network --- the node that, if removed, would most reduce temporal spread --- temporal
betweenness tells us it is C, not the high-static-degree nodes A or C (which happen to tie).
A static betweenness calculation might rank A higher because of its three static neighbours,
but temporal analysis correctly identifies C as the temporal hub.
\end{keyinsight}

%=============================================================================
\section{Inter-Event Time Distributions and Burstiness}
\label{sec:interevent_times}
%=============================================================================

\subsection{What is an Inter-Event Time?}

Consider a single node $v$ that participates in a series of contacts at times
$t_1 < t_2 < t_3 < \cdots < t_m$. The \textbf{inter-event times} are the gaps between
consecutive contacts:
\[
  \tau_1 = t_2 - t_1, \quad \tau_2 = t_3 - t_2, \quad \ldots, \quad \tau_{m-1} = t_m - t_{m-1}.
\]
Each $\tau_i$ is a positive number representing how long node $v$ waited between one contact
and the next.

\subsection{Poisson vs.\ Bursty Contact Patterns}

\begin{analogybox}{Bus Arrivals: Regular vs.\ Bursty}
Imagine waiting at a bus stop. In the ideal world, buses arrive exactly every 10 minutes:
perfectly \textbf{regular} (periodic). In a less ideal world, buses come at random but with an
average of one every 10 minutes: a \textbf{Poisson process}. In the real world, buses sometimes
cluster: three come in a row (a ``burst''), then nothing for 40 minutes.

The third pattern is called \textbf{bursty}: long silences punctuated by intense bursts of
activity. This is precisely what contact data looks like in real epidemiological settings.
Hospital staff talk intensely during morning rounds (a burst), then have long gaps during
administrative hours. The same pattern appears in phone calls, email, and face-to-face contact
sensors.

\textbf{Why does burstiness matter for epidemics?} Consider two nodes, each with an average
of 1 contact per hour:
\begin{itemize}[noitemsep]
  \item \textbf{Regular node}: contacts every hour, on the clock.
  \item \textbf{Bursty node}: 5 contacts within a 10-minute window, then silence for 5 hours.
\end{itemize}
The bursty node has long silent periods during which a pathogen it is carrying would cause
recovery before transmission. In the short burst windows, it transmits rapidly. The net effect
is often \textit{slower} epidemic spread than the regular node, despite the same average rate.
\end{analogybox}

\begin{figure}[H]
\centering
\begin{tikzpicture}[
  >=Stealth, thick,
  label/.style={font=\small\bfseries\sffamily}
]
% Left panel: regular (Poisson-like) contact sequence
\begin{scope}[xshift=0cm]
  \node[label] at (3, 2.5) {Regular (Poisson-like)};
  \draw[->] (0,0) -- (6.5,0) node[right,font=\scriptsize]{time};
  \draw[->] (0,0) -- (0,2.2) node[above,font=\scriptsize]{events};
  % Bars at roughly equal spacing
  \foreach \x in {0.5, 1.2, 1.9, 2.6, 3.4, 4.1, 4.8, 5.5}{
    \draw[blue!70!black, ultra thick] (\x,0) -- (\x,1.5);
    \filldraw[blue!70!black] (\x,0) circle (2pt);
  }
  \node[font=\scriptsize, color=blue!60!black] at (3, -0.4) {gaps $\approx$ equal};
\end{scope}

% Right panel: bursty contact sequence
\begin{scope}[xshift=8cm]
  \node[label] at (3, 2.5) {Bursty};
  \draw[->] (0,0) -- (6.5,0) node[right,font=\scriptsize]{time};
  \draw[->] (0,0) -- (0,2.2) node[above,font=\scriptsize]{events};
  % Burst 1: three close events
  \foreach \x in {0.3, 0.5, 0.7}{
    \draw[red!70!black, ultra thick] (\x,0) -- (\x,1.5);
    \filldraw[red!70!black] (\x,0) circle (2pt);
  }
  % Long gap then Burst 2
  \foreach \x in {4.0, 4.2, 4.4, 4.6}{
    \draw[red!70!black, ultra thick] (\x,0) -- (\x,1.5);
    \filldraw[red!70!black] (\x,0) circle (2pt);
  }
  % Annotation for gap
  \draw[<->, gray!60, thick] (0.7, -0.25) -- node[below,font=\scriptsize]{long silence} (4.0, -0.25);
  \node[font=\scriptsize, color=red!60!black] at (3, -0.6) {bursts + long gaps};
\end{scope}
\end{tikzpicture}
\caption{Comparison of a regular (Poisson-like) contact sequence (left) and a bursty contact
sequence (right). The bursty sequence has the same number of events but concentrates them in
short bursts separated by long silences. This pattern is ubiquitous in real contact data
(hospital wards, phone logs, proximity sensors) and has significant consequences for epidemic
dynamics.}
\label{fig:burstiness}
\end{figure}

\subsection{Quantifying Burstiness: The $B$ Parameter}

\begin{mathbox}{The Burstiness Parameter}
Let $\mu_\tau = \mathbb{E}[\tau]$ be the \textbf{mean} inter-event time and $\sigma_\tau =
\sqrt{\mathrm{Var}[\tau]}$ be the \textbf{standard deviation} of inter-event times for a node.
The \textbf{burstiness parameter} is:
\[
  B = \frac{\sigma_\tau - \mu_\tau}{\sigma_\tau + \mu_\tau} \in (-1,\, 1).
\]

\textbf{Interpreting the three limiting cases:}
\begin{itemize}
  \item \textbf{Perfectly periodic} ($\sigma_\tau = 0$): all gaps are equal, $\sigma_\tau = 0$.
        Then $B = (0 - \mu_\tau)/(0 + \mu_\tau) = -1$. Maximum negative burstiness: completely
        regular activity.
  \item \textbf{Poisson process} ($\tau \sim \text{Exp}(\lambda)$): for an exponential
        distribution, the mean equals the standard deviation, $\sigma_\tau = \mu_\tau$. Then
        $B = 0$. Neutral: random but not particularly bursty.
  \item \textbf{Maximally bursty} ($\sigma_\tau \gg \mu_\tau$): a few very large gaps dominate,
        so $\sigma_\tau \to \infty$ relative to $\mu_\tau$. Then $B \to 1$. Maximum positive
        burstiness: intense bursts separated by very long silences.
\end{itemize}

\textbf{Empirical finding:} Real human contact networks typically have $B \in [0.2, 0.8]$,
indicating moderately bursty activity. This is why epidemic simulation cannot simply assume
Poisson contacts.
\end{mathbox}

\begin{examplebox}{Computing Burstiness for a Node}
Suppose node $v$ has contacts at times $t = \{1, 2, 3, 10, 11, 12\}$. The inter-event times
are:
\[
  \tau = \{t_2 - t_1,\, t_3 - t_2,\, t_4 - t_3,\, t_5 - t_4,\, t_6 - t_5\}
       = \{1, 1, 7, 1, 1\}.
\]

\textbf{Mean:}
\[
  \mu_\tau = \frac{1 + 1 + 7 + 1 + 1}{5} = \frac{11}{5} = 2.2
\]

\textbf{Standard deviation:}
\[
  \sigma_\tau = \sqrt{\frac{(1-2.2)^2 + (1-2.2)^2 + (7-2.2)^2 + (1-2.2)^2 + (1-2.2)^2}{5}}
\]
\[
  = \sqrt{\frac{1.44 + 1.44 + 23.04 + 1.44 + 1.44}{5}}
  = \sqrt{\frac{28.8}{5}} = \sqrt{5.76} = 2.4
\]

\textbf{Burstiness:}
\[
  B = \frac{2.4 - 2.2}{2.4 + 2.2} = \frac{0.2}{4.6} \approx 0.043
\]

This is nearly Poisson ($B \approx 0$). Intuitively: although we see two ``bursts'' of
contacts, the single large gap ($\tau_3 = 7$) is only $3\times$ the mean gap. True bursty
sequences have gaps that are $10$--$100\times$ the mean.

\textbf{Compare with a truly bursty node} whose inter-event times are $\{0.1, 0.1, 0.1, 50,
0.1\}$:
\[
  \mu_\tau = \frac{50.4}{5} = 10.08, \quad
  \sigma_\tau \approx 19.9, \quad
  B = \frac{19.9 - 10.08}{19.9 + 10.08} \approx 0.33.
\]
This is moderately bursty ($B \approx 0.33$), reflecting the one long gap of 50 time units
surrounded by rapid-fire contacts.
\end{examplebox}

\begin{keyinsight}{Burstiness and Source Detection}
A node that is bursty may look less infectious than its average contact rate suggests: its
long silent periods give the pathogen time to ``time out'' (the $\Delta t$ constraint). When
fitting epidemic models or choosing the $\Delta t$ parameter in a $\Delta t$-constrained path
analysis, bursty contact distributions require special attention. Incorrectly assuming Poisson
contacts in a bursty network will lead to wrong reachability estimates and, consequently, wrong
source predictions.
\end{keyinsight}

%=============================================================================
\section{How the Codebase Represents Temporal Networks}
\label{sec:codebase_representation}
%=============================================================================

Having built the theory from scratch, we now close the loop by showing exactly how these
concepts are implemented in the project codebase. Every abstract object introduced in this
chapter has a concrete counterpart in the code.

\subsection{The Event List Format}

The project stores temporal networks as a list of $(u, v, t)$ triples, sorted by time. This is
the COO format described in Section~\ref{sec:adjacency_representation}. In Python, after
loading via NetworkX, the function \texttt{make\_array\_from\_networkx} (in
\texttt{setup/read\_network.py}) converts the graph object to a plain NumPy array:

\begin{codewalk}{make\_array\_from\_networkx in setup/read\_network.py}
\begin{lstlisting}[language=Python]
def make_array_from_networkx(H):
    """From the networkx graph, make a numpy array
    with rows (u, v, t), sorted by t."""
    rows = []
    for u, v, data in H.edges(data=True):
        # Each edge stores a list of timestamps
        rows.extend([[u, v, t] for t in data['times']])
    arr = np.array(rows)
    arr = arr[arr[:, 2].argsort()]  # sort by time
    return arr
\end{lstlisting}

The result is a NumPy array with shape \texttt{(|E\^{}T|, 3)}, where:
\begin{itemize}[noitemsep]
  \item Column 0: source node $u$ (integer index)
  \item Column 1: target node $v$ (integer index)
  \item Column 2: timestamp $t$
\end{itemize}

For an Alice--Bob--Carol style example, this would look like:
\begin{lstlisting}[language=Python]
# arr for a 3-node, 3-event network:
arr = np.array([
    [0, 1, 1],   # Alice(0) -- Bob(1)   at t=1
    [1, 2, 2],   # Bob(1)   -- Carol(2) at t=2
    [2, 0, 7],   # Carol(2) -- Alice(0) at t=7
])
\end{lstlisting}
\end{codewalk}

\subsection{Loading the Network}

The entry point for loading a temporal network is \texttt{setup/read\_network.py}. It handles
two cases:

\begin{codewalk}{load\_network in setup/read\_network.py}
\begin{lstlisting}[language=Python]
def load_network(cfg, label_attribute="old_id"):
    if cfg.nwk.type == "empirical":
        # Load from CSV file: each row is "u v t"
        H = read_networkx(
            'nwk/' + cfg.nwk.name + '.csv',
            t_max=cfg.nwk.t_max,
            directed=cfg.nwk.directed
        )
    elif cfg.nwk.type == "synthetic":
        # Generate e.g. Erdos-Renyi with all time steps
        H = generate_synthetic_graph(cfg)
    return H
\end{lstlisting}

For \textbf{empirical networks}, each line of the CSV file contains \texttt{u v t}: the two
node labels and the timestamp. Events after \texttt{t\_max} are ignored. Self-contacts
(where $u = v$) are discarded. The function returns a NetworkX graph object where each edge
\texttt{(u, v)} carries a \texttt{'times'} attribute: a list of all timestamps at which that
pair was in contact.

For \textbf{synthetic networks} (e.g.\ Erd\H{o}s--R\'enyi graphs), the contact times are set
to \texttt{range(t\_max + 1)}: every edge is active at every discrete time step. This is the
``always-on'' baseline that is equivalent to a static graph for temporal methods.
\end{codewalk}

\subsection{The TSIRData Dataclass}

Once simulations have been run, all data --- network, simulation outcomes, and ground truth ---
is stored in a single \texttt{TSIRData} dataclass (in \texttt{setup/data\_loader.py}):

\begin{codewalk}{TSIRData Dataclass in setup/data\_loader.py}
\begin{lstlisting}[language=Python]
@dataclass
class TSIRData:
    config: dict         # experiment configuration
    n_nodes: int         # |V|: number of nodes
    n_runs: int          # ground-truth simulation runs
    mc_runs: int         # Monte Carlo simulation runs
    # Monte Carlo SIR state arrays, shape: (n_nodes, mc_runs, n_nodes)
    mc_S: np.ndarray     # mc_S[s, r, v] = 1 if node v is Susceptible
    mc_I: np.ndarray     #   in run r of MC sim starting at source s
    mc_R: np.ndarray
    # Ground truth SIR state arrays, shape: (n_nodes, n_runs, n_nodes)
    truth_S: np.ndarray  # truth_S[s, r, v] = 1 if node v is Susceptible
    truth_I: np.ndarray  #   in ground-truth run r with true source s
    truth_R: np.ndarray
    # Binary masks for filtering:
    maximal_outbreak: np.ndarray  # shape (n_nodes, n_nodes)
    possible: np.ndarray          # shape (n_nodes, n_runs, n_nodes)
    lik_possible: np.ndarray      # log-likelihood correction mask
\end{lstlisting}

The three-dimensional arrays deserve explanation:
\begin{itemize}
  \item \texttt{truth\_S[s, r, v]}: the SIR state of node $v$ in run $r$ of the ground-truth
        simulation, given that node $s$ was the true source. This is the label data our GNN
        trains on.
  \item \texttt{mc\_S[s, r, v]}: the SIR state of node $v$ in Monte Carlo run $r$ with
        seed node $s$. These are the simulated ``observations'' the model uses for inference.
  \item \texttt{maximal\_outbreak[s, v]}: 1 if node $v$ \textit{ever} becomes infected in any
        run with source $s$, 0 otherwise. This is exactly the temporal reachability indicator
        $\mathbf{1}[s \leadsto v]$ --- a direct implementation of the temporal reachability
        concept from Section~\ref{sec:reachability_comparison}.
\end{itemize}
\end{codewalk}

\begin{keyinsight}{Theory Meets Code}
Every concept in this chapter has a direct counterpart in the codebase:
\begin{itemize}[noitemsep]
  \item The event list $(u, v, t)$ $\Leftrightarrow$ the \texttt{arr} array from
        \texttt{make\_array\_from\_networkx}.
  \item The observation window $\mathcal{T}$ $\Leftrightarrow$ \texttt{cfg.nwk.t\_max}.
  \item The $\Delta t$ constraint $\Leftrightarrow$ the SIR parameter \texttt{cfg.sir.mu}
        (recovery rate; $\Delta t \approx 1/\mu$).
  \item Temporal reachability $\mathcal{I}(s)$ $\Leftrightarrow$
        \texttt{TSIRData.maximal\_outbreak[s, :]}.
  \item The static projection $G$ $\Leftrightarrow$ the NetworkX graph $H$ (which stores the
        edge list but ignores timing when passed to a static GNN).
\end{itemize}
The entire research project is, at its core, an attempt to exploit the temporal information
that the static projection discards --- to use $E^T$ instead of $G$ whenever possible.
\end{keyinsight}

%=============================================================================
\section{Static Projection and Information Loss}
\label{sec:static_projection}
%=============================================================================

We formalise what we lose when we project a temporal network to a static graph.

\begin{mathbox}{The Static Projection: Formal Definition}
Given a temporal network $G^T = (\mathcal{V}, E^T, \mathcal{T})$, its \textbf{static
projection} is the undirected graph $G = (\mathcal{V}, \mathcal{E})$ where:
\[
  (u, v) \in \mathcal{E} \iff \exists\, t \in \mathcal{T} : (u, v, t) \in E^T.
\]
In words: two nodes are connected in $G$ if and only if they were in contact at
\textit{any} point during the observation window.
\end{mathbox}

\begin{figure}[H]
\centering
\begin{tikzpicture}[
    every node/.style={font=\small},
    event/.style={circle, draw, thick, fill=white, minimum size=8mm,
                  font=\small\bfseries},
    timelabel/.style={font=\footnotesize, color=gray!70!black},
    evtarrow/.style={-{Stealth[length=2.5mm]}, thick},
    tstamp/.style={above, font=\scriptsize, color=blue!60!black},
    projlabel/.style={font=\small\sffamily\bfseries},
]

% -------- LEFT: Temporal Network --------
\begin{scope}[xshift=0cm]
  \node[projlabel] at (3.0, 5.0) {Temporal Contact Sequence $G^T$};

  % Time axis
  \draw[-{Stealth}, thick, gray!60] (-0.3, -0.2) -- (6.3, -0.2)
        node[right, font=\footnotesize]{$t$};
  \foreach \x/\lab in {0/1, 1.5/2, 3.0/3, 4.5/4, 6.0/5}{
    \draw[gray!60] (\x, -0.1) -- (\x, -0.3);
    \node[timelabel] at (\x, -0.55) {\lab};
  }

  % Nodes on y-axis
  \foreach \y/\lbl in {0.6/D, 1.6/C, 2.6/B, 3.6/A}{
    \node[left, font=\small\bfseries] at (-0.3, \y) {\lbl};
    \draw[dotted, gray!30] (-0.2, \y) -- (6.2, \y);
  }

  % Events as vertical bars
  % e1: A-B at t=1
  \draw[blue!70!black, ultra thick] (0, 2.6) -- (0, 3.6);
  \filldraw[blue!70!black] (0,2.6) circle (2.5pt);
  \filldraw[blue!70!black] (0,3.6) circle (2.5pt);
  \node[tstamp, color=blue!70!black] at (0.0, 3.8) {$e_1,t{=}1$};

  % e2: B-C at t=2
  \draw[blue!70!black, ultra thick] (1.5, 1.6) -- (1.5, 2.6);
  \filldraw[blue!70!black] (1.5,1.6) circle (2.5pt);
  \filldraw[blue!70!black] (1.5,2.6) circle (2.5pt);
  \node[tstamp, color=blue!70!black] at (1.5, 2.8) {$e_2,t{=}2$};

  % e3: C-D at t=3
  \draw[gray!60!black, thick] (3.0, 0.6) -- (3.0, 1.6);
  \filldraw[gray!60!black] (3.0,0.6) circle (2.5pt);
  \filldraw[gray!60!black] (3.0,1.6) circle (2.5pt);
  \node[tstamp, color=gray!50!black] at (3.0, 1.8) {$e_3,t{=}3$};

  % e4: A-C at t=4  (this is the "phantom builder" in static graph)
  \draw[orange!80!black, ultra thick] (4.5, 1.6) -- (4.5, 3.6);
  \filldraw[orange!80!black] (4.5,1.6) circle (2.5pt);
  \filldraw[orange!80!black] (4.5,3.6) circle (2.5pt);
  \node[tstamp, color=orange!70!black] at (4.5, 3.8) {$e_4,t{=}4$};

  % e5: D-A at t=5
  \draw[gray!60!black, thick] (6.0, 0.6) -- (6.0, 3.6);
  \filldraw[gray!60!black] (6.0,0.6) circle (2.5pt);
  \filldraw[gray!60!black] (6.0,3.6) circle (2.5pt);
  \node[tstamp, color=gray!50!black] at (6.0, 3.8) {$e_5,t{=}5$};

  % Phantom path annotation
  \node[red!70!black, font=\scriptsize, align=center] at (3.0, -1.2)
    {\textbf{Static graph says:} $C \to B \to A$ exists.\\
     \textbf{Temporal graph: impossible.}\\
     $B$'s last event is $t=2$; $C$'s first outgoing is $t=3>2$.};
\end{scope}

% -------- RIGHT: Static Projection --------
\begin{scope}[xshift=9.5cm, yshift=0.5cm]
  \node[projlabel] at (1.0, 3.8) {Static Projection $G$};

  \node[event, fill=blue!8]   (A) at (0.0, 3.0) {A};
  \node[event, fill=blue!8]   (B) at (0.0, 1.5) {B};
  \node[event, fill=blue!8]   (C) at (2.0, 1.5) {C};
  \node[event, fill=blue!8]   (D) at (2.0, 3.0) {D};

  \draw[thick] (A)--(B);
  \draw[thick] (B)--(C);
  \draw[thick] (C)--(D);
  \draw[thick] (A)--(C);
  \draw[thick] (D)--(A);

  % Highlight the phantom path C->B->A
  \draw[red!70!black, ultra thick, dashed, ->, bend left=30]
    (C) to (B);
  \draw[red!70!black, ultra thick, dashed, ->, bend left=30]
    (B) to (A);

  \node[red!70!black, font=\scriptsize, align=center] at (1.0, -0.4)
    {Path $C{\to}B{\to}A$: exists in $G$\\but is a \textbf{phantom} in $G^T$};
\end{scope}

\end{tikzpicture}
\caption{Left: a temporal contact sequence on four nodes with five events. Nodes are arranged
on the $y$-axis; time runs left to right. The blue and orange events form valid forward paths.
Right: the static projection retains all edges but loses timestamps, creating the phantom path
$C \to B \to A$ (dashed red) which is causally impossible because B's last event ($t=2$)
precedes C's first outgoing event ($t=3$).}
\label{fig:static_projection}
\end{figure}

\begin{keyinsight}{Quantifying Information Loss}
Let $\mathcal{R}^T(u)$ denote the \textbf{temporal reachability set} of node $u$ and
$\mathcal{R}(u)$ its static reachability set. By construction $\mathcal{R}^T(u) \subseteq
\mathcal{R}(u)$ always holds. In empirical contact networks with strong burstiness, the ratio
$|\mathcal{R}^T(u)| / |\mathcal{R}(u)|$ can be as low as $0.1$--$0.3$: the static graph
over-estimates causal reach by a factor of 3--10. For source detection, this means the static
graph floods the posterior over sources with false candidates; a static GNN may assign high
probability to nodes that could never have initiated the observed outbreak.
\end{keyinsight}

\begin{analogybox}{Collapsing a Train Timetable to a Route Map}
A train timetable tells you that Train 42 leaves Munich at 08:15 and arrives Augsburg at 08:52,
while Train 7 leaves Augsburg at 08:40 for Nuremberg. The route map shows
Munich--Augsburg--Nuremberg as a connected path. But the timetable reveals the connection is
impossible: you arrive at 08:52 and your connection left at 08:40.

Collapsing a temporal network into a static graph is exactly like reducing the timetable to a
route map: you preserve topology but destroy the information needed to reason about what is
actually reachable.
\end{analogybox}

%=============================================================================
\section{The Temporal Event Graph: A Forward Pointer}
\label{sec:event_graph_pointer}
%=============================================================================

All of the path and reachability concepts developed in this chapter reach their most
computationally powerful form in the \textbf{Temporal Event Graph (TEG)} representation due
to Saram\"{a}ki, Kivel\"{a}, and Karsai. Rather than treating nodes as primary objects, the
TEG elevates \textit{events} to nodes in a static directed acyclic graph (DAG), where directed
edges encode the possibility of transmission between consecutive contacts.

\begin{keyinsight}{Why the TEG Matters for GNN-Based Source Detection}
Every concept in this chapter maps cleanly onto the TEG structure:
\begin{itemize}[noitemsep]
  \item \textbf{Time-respecting paths} in $G^T$ become \textbf{directed paths} in the TEG.
  \item \textbf{$\Delta t$-constrained paths} correspond to edges with inter-event gap
        $\leq \Delta t$ in the TEG.
  \item \textbf{Temporal reachability} corresponds to graph reachability in the DAG.
  \item \textbf{Asymmetric reachability} is automatic in a DAG (acyclicity guarantees it).
\end{itemize}
A GNN operating on the TEG therefore inherits the full causal structure of the temporal network
without any approximation. This is the theoretical justification for the DAG-GNN architecture
described in Chapter~10.
\end{keyinsight}

Chapter~10 covers the full TEG construction algorithm, its weighted variant (edge weights encode
inter-event gaps), and the DAG convolutional architecture that performs backward message passing
over it. The key forward pointer: when our DAG-GNN flips the TEG edges and runs message passing
toward earlier events, it is precisely performing the temporal reachability computation in
reverse --- tracing the epidemic backward toward its origin.

%=============================================================================
\section{Summary}
\label{sec:tn_summary}
%=============================================================================

This chapter has built the complete formal toolkit for working with temporal networks. Here is
a concise summary of every concept introduced, in the order they appeared:

\begin{conceptbox}{Chapter Summary}
\begin{itemize}
  \item A \textbf{temporal network} $G^T = (\mathcal{V}, E^T, \mathcal{T})$ records contacts
        as timestamped triples $(u, v, t)$. The \textbf{static projection} discards all
        timestamps and retains only connectivity.
  \item \textbf{Phantom paths} are paths that exist in the static projection but violate
        causal time-ordering. They are the fundamental reason static GNNs produce incorrect
        source posterior estimates.
  \item A \textbf{time-respecting path} requires strictly increasing event times:
        $t_1 < t_2 < \cdots < t_k$. This is the only type of path that can carry a pathogen.
  \item A \textbf{$\Delta t$-constrained path} additionally requires that no gap between
        consecutive events exceeds $\Delta t$, the maximum infectious duration. Smaller
        $\Delta t$ prunes more paths and makes source detection harder.
  \item \textbf{Temporal reachability} is \textit{asymmetric}: $u \leadsto v$ does not imply
        $v \leadsto u$, even in undirected networks.
  \item The \textbf{adjacency tensor} $A_{ij}(t)$ captures temporal contacts in three
        dimensions. In practice, we store it as a COO event list, the direct input to
        PyTorch Geometric's \texttt{edge\_index}.
  \item \textbf{Temporal degree} counts total contact events per node.
        \textbf{Temporal betweenness} counts how often a node lies on shortest
        time-respecting paths. Both can differ substantially from their static counterparts.
  \item \textbf{Burstiness} $B \in (-1, 1)$ quantifies whether a node's contacts are regular
        ($B < 0$), random ($B \approx 0$), or clustered into bursts ($B > 0$). Real contact
        networks are typically bursty, which slows epidemic spread and complicates inference.
  \item In the codebase, the event list lives in the NetworkX graph \texttt{H} (loaded by
        \texttt{setup/read\_network.py}), the temporal reachability is encoded in
        \texttt{TSIRData.maximal\_outbreak}, and the observation window is controlled by
        \texttt{cfg.nwk.t\_max}.
\end{itemize}
\end{conceptbox}

\begin{keyinsight}{The Central Lesson for Source Detection}
Every simplification made when moving from a temporal contact sequence to a static graph
introduces \textit{false causal paths}: connections the static model treats as evidence for or
against a candidate source, but which could never have been traversed by the actual pathogen.
The further we are from the original event sequence, the noisier the posterior over sources
becomes. Chapters~4--10 develop model architectures that progressively close this gap, from
static GCN baselines through to DAG-GNNs that operate directly on the causal event structure.
\end{keyinsight}
